% Source: https://artofproblemsolving.com/community/c3h1447123_100th_post_special_marathon
% Stable identity: aops-topic-1447123-q01
Each of the $20$ triangular faces has $3$ edges, which gives $20\cdot3=60$ edge-face incidences. Every edge of a polyhedron lies on exactly two faces, so each edge has been counted twice:
\[
E=\dfrac{20\cdot3}{2}=30 .
\]

\begin{asy}
usepackage("amsmath");
size(5cm);
defaultpen(fontsize(10pt));
pair[] V={(0.9519,-1.6071),(1.8668,0.0175),(0.5813,-0.9999),(0.9192,1.6287),(-1.1608,-0.0175),(-0.5611,-0.9999),(-0.9192,-1.6287),(1.1608,0.0175),(0.5611,0.9999),(-0.9519,1.6071),(-1.8668,-0.0175),(-0.5813,0.9999)};
draw(V[0]--V[2],gray+dashed);
draw(V[1]--V[2],gray+dashed);
draw(V[1]--V[8],gray+dashed);
draw(V[2]--V[4],gray+dashed);
draw(V[2]--V[6],gray+dashed);
draw(V[2]--V[8],gray+dashed);
draw(V[3]--V[8],red+dashed);
draw(V[4]--V[6],gray+dashed);
draw(V[4]--V[8],gray+dashed);
draw(V[4]--V[9],gray+dashed);
draw(V[4]--V[10],gray+dashed);
draw(V[8]--V[9],gray+dashed);
draw(V[0]--V[1],linewidth(0.8));
draw(V[0]--V[5],linewidth(0.8));
draw(V[0]--V[6],linewidth(0.8));
draw(V[0]--V[7],linewidth(0.8));
draw(V[1]--V[3],red+linewidth(1.1));
draw(V[1]--V[7],linewidth(0.8));
draw(V[3]--V[7],red+linewidth(1.1));
draw(V[3]--V[9],red+linewidth(1.1));
draw(V[3]--V[11],red+linewidth(1.1));
draw(V[5]--V[6],linewidth(0.8));
draw(V[5]--V[7],linewidth(0.8));
draw(V[5]--V[10],linewidth(0.8));
draw(V[5]--V[11],linewidth(0.8));
draw(V[6]--V[10],linewidth(0.8));
draw(V[7]--V[11],linewidth(0.8));
draw(V[9]--V[10],linewidth(0.8));
draw(V[9]--V[11],linewidth(0.8));
draw(V[10]--V[11],linewidth(0.8));
for(int i=0;i<12;++i) dot(V[i],black+3);
dot(V[3],red+5);
label("\mbox{}",(min(currentpicture,user=true)+max(currentpicture,user=true))/2,black); // keeps the text after the figure black when LaTeX typesets the labels inline
\end{asy}

(As a check, five faces meet at each vertex, so the number of vertices is $V=\dfrac{20\cdot3}{5}=12$, and $V-E+F=12-30+20=2$, as Euler's formula requires. In the figure hidden edges are dashed, and the five edges at the large red vertex are red; one of them is hidden.)

The answer is $\boxed{\textbf{(A) } 30}$.
